% GENERATED FILE. Edit canonical lessons, then run npm run generate:books.
% canonical-source-hash: fnv1a64:277280042ff31ea4
% canonical-lessons: TA-C458-kataru, TA-C458-vimmu, TA-C458-vicumpu, TA-C458-pulampu, TA-C458-alaru

\chapter{Crying Aloud}
\label{ch:ta-crying-aloud}
\hlchaptermodality{458}

\begin{chapteropening}
\textbf{By the end of this chapter:} \emph{I can describe ways of crying, sobbing and screaming.}

Complete the last lesson of chapter 458: I can describe ways of crying, sobbing and screaming.
\end{chapteropening}

\section[\texorpdfstring{\ta{கதறு} (kataṟu)}{கதறு (kataṟu)}]{\ta{கதறு} (kataṟu) — to wail; to cry out}

\label{lesson:TA-C458-kataru}



\begin{quote}
Before the new one: say the Tamil for pitifulness, then the Tamil for concern.
\end{quote}

\subsection*{You'll want to know: \ta{கதறு}}
\textbf{\ta{கதறு}} — \emph{kataṟu} — \textquotedblleft{}to wail; to cry out\textquotedblright{}.

\begin{quote}
\ta{அந்தச்} \ta{செய்தியைக்} \ta{கேட்டு} \ta{அவள்} \ta{கதறி} \ta{அழுதாள்}.
\end{quote}

\begin{quote}
\emph{antac ceytiyaik kēṭṭu avaḷ kataṟi aḻutāḷ.}
\end{quote}

\begin{quote}
Hearing the news, she wailed and wept.
\end{quote}

Past (I …): \ta{கதறினேன்}, \emph{kataṟiṉēṉ}.

Often \ta{கதறி} \ta{அழு}, to cry aloud.

\begin{cousinweb}
\textbf{\ta{கதறு}} (\emph{kataṟu}) is native Dravidian (DEDR 1189); Malayalam has the same verb for roaring.
\end{cousinweb}

\begin{grammarlens}[title={What you've built}]
One more word for describing what you lived and felt.
\end{grammarlens}

\subsection*{Your turn}
\begin{itemize}
\raggedright
  \item \emph{Say it:} \emph{kataṟu}
  \item \emph{Say it:} \emph{kataṟu}, once more
  \item \emph{Say it:} \emph{antac ceytiyaik kēṭṭu avaḷ kataṟi aḻutāḷ.} — Hearing the news, she wailed and wept.
  \item \emph{Recall:} say the Tamil for pitifulness, then the Tamil for concern, then say \emph{kataṟu} again
\end{itemize}

\begin{tcolorbox}[breakable,colback=teal!4,colframe=teal!35!black,title={Before you move on}]
What does \ta{கதறு} mean? (\textquotedblleft{}To wail; to cry out\textquotedblright{}.) Say it once more.
\end{tcolorbox}
\section[\texorpdfstring{\ta{விம்மு} (vimmu)}{விம்மு (vimmu)}]{\ta{விம்மு} (vimmu) — to sob}

\label{lesson:TA-C458-vimmu}



\begin{quote}
Before the new one: say the Tamil for concern, then the Tamil for to wail.
\end{quote}

\subsection*{You'll want to know: \ta{விம்மு}}
\textbf{\ta{விம்மு}} — \emph{vimmu} — \textquotedblleft{}to sob\textquotedblright{}.

\begin{quote}
\ta{குழந்தை} \ta{அம்மாவைத்} \ta{தேடி} \ta{விம்மியது}.
\end{quote}

\begin{quote}
\emph{kuḻantai ammāvait tēṭi vimmiyatu.}
\end{quote}

\begin{quote}
The child sobbed, looking for its mother.
\end{quote}

Past (I …): \ta{விம்மினேன்}, \emph{vimmiṉēṉ}.

Past with -\ta{இன்}-: \ta{விம்மினேன்}. \ta{விம்மல்}, a sob.

\begin{cousinweb}
\textbf{\ta{விம்மு}} (\emph{vimmu}) is native Dravidian (DEDR 5402). DEDR keeps a second \ta{விம்மு}, \textquotedblleft{}to swell\textquotedblright{}, apart as its own entry (5397).
\end{cousinweb}

\begin{grammarlens}[title={What you've built}]
One more word for describing what you lived and felt.
\end{grammarlens}

\subsection*{Your turn}
\begin{itemize}
\raggedright
  \item \emph{Say it:} \emph{vimmu}
  \item \emph{Say it:} \emph{vimmu}, once more
  \item \emph{Say it:} \emph{kuḻantai ammāvait tēṭi vimmiyatu.} — The child sobbed, looking for its mother.
  \item \emph{Recall:} say the Tamil for concern, then the Tamil for to wail, then say \emph{vimmu} again
\end{itemize}

\begin{tcolorbox}[breakable,colback=teal!4,colframe=teal!35!black,title={Before you move on}]
What does \ta{விம்மு} mean? (\textquotedblleft{}To sob\textquotedblright{}.) Say it once more.
\end{tcolorbox}
\section[\texorpdfstring{\ta{விசும்பு} (vicumpu)}{விசும்பு (vicumpu)}]{\ta{விசும்பு} (vicumpu) — to whimper; to sob softly}

\label{lesson:TA-C458-vicumpu}



\begin{quote}
Before the new one: say the Tamil for to wail, then the Tamil for to sob.
\end{quote}

\subsection*{You'll want to know: \ta{விசும்பு}}
\textbf{\ta{விசும்பு}} — \emph{vicumpu} — \textquotedblleft{}to whimper; to sob softly\textquotedblright{}.

\begin{quote}
\ta{தோற்ற} \ta{பிறகு} \ta{அவன்} \ta{தனியாக} \ta{விசும்பினான்}.
\end{quote}

\begin{quote}
\emph{tōṟṟa piṟaku avaṉ taṉiyāka vicumpiṉāṉ.}
\end{quote}

\begin{quote}
After losing, he whimpered on his own.
\end{quote}

Past (I …): \ta{விசும்பினேன்}, \emph{vicumpiṉēṉ}.

Often doubled: \ta{விசும்பி} \ta{விசும்பி} \ta{அழுதாள்}, she sobbed and sobbed.

\begin{cousinweb}
\textbf{\ta{விசும்பு}} (\emph{vicumpu}) is a Tamil verb whose origin is not settled; DEDR has no entry for this sense.
\end{cousinweb}

\begin{grammarlens}[title={What you've built}]
One more word for describing what you lived and felt.
\end{grammarlens}

\subsection*{Your turn}
\begin{itemize}
\raggedright
  \item \emph{Say it:} \emph{vicumpu}
  \item \emph{Say it:} \emph{vicumpu}, once more
  \item \emph{Say it:} \emph{tōṟṟa piṟaku avaṉ taṉiyāka vicumpiṉāṉ.} — After losing, he whimpered on his own.
  \item \emph{Recall:} say the Tamil for to wail, then the Tamil for to sob, then say \emph{vicumpu} again
\end{itemize}

\begin{tcolorbox}[breakable,colback=teal!4,colframe=teal!35!black,title={Before you move on}]
What does \ta{விசும்பு} mean? (\textquotedblleft{}To whimper; to sob softly\textquotedblright{}.) Say it once more.
\end{tcolorbox}
\section[\texorpdfstring{\ta{புலம்பு} (pulampu)}{புலம்பு (pulampu)}]{\ta{புலம்பு} (pulampu) — to lament; to moan}

\label{lesson:TA-C458-pulampu}



\begin{quote}
Before the new one: say the Tamil for to sob, then the Tamil for to whimper.
\end{quote}

\subsection*{You'll want to know: \ta{புலம்பு}}
\textbf{\ta{புலம்பு}} — \emph{pulampu} — \textquotedblleft{}to lament; to moan\textquotedblright{}.

\begin{quote}
\ta{அவர்} \ta{இரவு} \ta{முழுவதும்} \ta{புலம்பினார்}.
\end{quote}

\begin{quote}
\emph{avar iravu muḻuvatum pulampiṉār.}
\end{quote}

\begin{quote}
He moaned all night long.
\end{quote}

Past (I …): \ta{புலம்பினேன்}, \emph{pulampiṉēṉ}.

Also of a feverish person talking aloud without sense.

\begin{cousinweb}
\textbf{\ta{புலம்பு}} (\emph{pulampu}) is native Dravidian (DEDR 4304); the entry runs from speaking foolishly to wailing in grief.
\end{cousinweb}

\begin{grammarlens}[title={What you've built}]
One more word for describing what you lived and felt.
\end{grammarlens}

\subsection*{Your turn}
\begin{itemize}
\raggedright
  \item \emph{Say it:} \emph{pulampu}
  \item \emph{Say it:} \emph{pulampu}, once more
  \item \emph{Say it:} \emph{avar iravu muḻuvatum pulampiṉār.} — He moaned all night long.
  \item \emph{Recall:} say the Tamil for to sob, then the Tamil for to whimper, then say \emph{pulampu} again
\end{itemize}

\begin{tcolorbox}[breakable,colback=teal!4,colframe=teal!35!black,title={Before you move on}]
What does \ta{புலம்பு} mean? (\textquotedblleft{}To lament; to moan\textquotedblright{}.) Say it once more.
\end{tcolorbox}
\section[\texorpdfstring{\ta{அலறு} (alaṟu)}{அலறு (alaṟu)}]{\ta{அலறு} (alaṟu) — to scream; to shriek}

\label{lesson:TA-C458-alaru}



\begin{quote}
Before the new one: say the Tamil for to whimper, then the Tamil for to lament.
\end{quote}

\subsection*{You'll want to know: \ta{அலறு}}
\textbf{\ta{அலறு}} — \emph{alaṟu} — \textquotedblleft{}to scream; to shriek\textquotedblright{}.

\begin{quote}
\ta{பாம்பைப்} \ta{பார்த்ததும்} \ta{அவள்} \ta{அலறினாள்}.
\end{quote}

\begin{quote}
\emph{pāmpaip pārttatum avaḷ alaṟiṉāḷ.}
\end{quote}

\begin{quote}
She screamed when she saw the snake.
\end{quote}

Past (I …): \ta{அலறினேன்}, \emph{alaṟiṉēṉ}.

The verbal participle: \ta{அலறி} \ta{ஓடினான்}, he ran off screaming.

\begin{cousinweb}
\textbf{\ta{அலறு}} (\emph{alaṟu}) is native Dravidian (DEDR 245), in an entry of words for noisy talk and loud cries.
\end{cousinweb}

\begin{grammarlens}[title={What you've built}]
One more word for describing what you lived and felt.
\end{grammarlens}

\subsection*{Your turn}
\begin{itemize}
\raggedright
  \item \emph{Say it:} \emph{alaṟu}
  \item \emph{Say it:} \emph{alaṟu}, once more
  \item \emph{Say it:} \emph{pāmpaip pārttatum avaḷ alaṟiṉāḷ.} — She screamed when she saw the snake.
  \item \emph{Recall:} say the Tamil for to whimper, then the Tamil for to lament, then say \emph{alaṟu} again
\end{itemize}

\begin{tcolorbox}[breakable,colback=teal!4,colframe=teal!35!black,title={Before you move on}]
What does \ta{அலறு} mean? (\textquotedblleft{}To scream; to shriek\textquotedblright{}.) Say it once more.
\end{tcolorbox}
